%% ============================================================================
%% Section: Quantitative Layer and Expressiveness
%% Included from main.tex via %% ============================================================================
%% Section: Quantitative Layer and Expressiveness
%% Included from main.tex via %% ============================================================================
%% Section: Quantitative Layer and Expressiveness
%% Included from main.tex via %% ============================================================================
%% Section: Quantitative Layer and Expressiveness
%% Included from main.tex via \input{sections/quantitative}
%% Faithful to LinearGaussianInference.java, CounterfactualInference.java,
%% DefaultCausalGraphRuntime.java. KEEP SHORT -- expressiveness, not a contribution.
%% ============================================================================
\section{Quantitative Layer and Expressiveness}\label{sec:quant}

The same compiled snapshot that answers structural queries also carries optional
numeric parameters, which let Causal4j answer \emph{quantitative} causal
questions---interventions and counterfactuals---without leaving the DBMS. We
describe the layer briefly: it is a demonstration of the expressiveness the
indexed engine enables, not a contribution of the paper.

\paragraph{Linear-Gaussian model}
When a model is parameterized, each variable follows the structural equation
$V = \mu_{0,V} + \sum_{P\in\mathrm{pa}(V)} \beta_{P\to V}\,P + \varepsilon_V$ with
$\varepsilon_V\sim\mathcal{N}(0,\sigma^2_V)$, where the edge coefficients
$\beta$, intercepts $\mu_0$, and noise variances $\sigma^2$ are exactly the
\textsf{weight}, \textsf{intercept}, and \textsf{noise} properties read at compile
time (\S\ref{sec:compile}). The snapshot reports itself parameterized only when
some value departs from the structural-only defaults, so a purely qualitative DAG
incurs no numeric machinery.

\paragraph{Interventions and counterfactuals}
The do-operator $do(X{=}x)$ fixes $X$'s mean to $x$, sets its variance to zero,
and drops its incoming edges; the general recursion then handles it with no
special case. Interventional means propagate in topological order,
$\mu(V)=\mu_{0,V}+\sum_P \beta_{P\to V}\,\mu(P)$; interventional variances follow
from an incremental covariance recursion in the same order; and the total effect
of $T$ on $Y$ is obtained by unit-impulse propagation
($\partial T{=}1$, $\partial V=\sum_P \beta_{P\to V}\,\partial P$, read off at
$Y$), avoiding any path enumeration. Counterfactuals follow the standard
abduction--action--prediction recipe over the same equations: infer the exogenous
noise consistent with the observed factual world, apply the hypothesized
intervention, and re-propagate. Every one of these reuses the topological order
already on the snapshot, so the quantitative layer adds semantics, not a second
data structure.

\paragraph{Worked example}
On a parameterized version of the running example, \lstinline[style=causaldsl]|INTERVENE ON interventions = 1 COMPARE TO interventions = 0 RETURN EFFECT ON reported_cases|
returns the total effect along the front-door paths, while a
\lstinline[style=causaldsl]|COUNTERFACTUAL| statement answers what
\textsf{reported\_cases} would have been for a specific observed unit under a
different intervention. Because $do(\cdot)$ yields a derived model---the DAG with
$X$'s in-edges removed---these statements register a new causal world in the
model registry (\S\ref{sec:lang:registry}) that subsequent structural queries can
address like any other, closing the loop between the qualitative and quantitative
layers over a single stored graph.

%% Faithful to LinearGaussianInference.java, CounterfactualInference.java,
%% DefaultCausalGraphRuntime.java. KEEP SHORT -- expressiveness, not a contribution.
%% ============================================================================
\section{Quantitative Layer and Expressiveness}\label{sec:quant}

The same compiled snapshot that answers structural queries also carries optional
numeric parameters, which let Causal4j answer \emph{quantitative} causal
questions---interventions and counterfactuals---without leaving the DBMS. We
describe the layer briefly: it is a demonstration of the expressiveness the
indexed engine enables, not a contribution of the paper.

\paragraph{Linear-Gaussian model}
When a model is parameterized, each variable follows the structural equation
$V = \mu_{0,V} + \sum_{P\in\mathrm{pa}(V)} \beta_{P\to V}\,P + \varepsilon_V$ with
$\varepsilon_V\sim\mathcal{N}(0,\sigma^2_V)$, where the edge coefficients
$\beta$, intercepts $\mu_0$, and noise variances $\sigma^2$ are exactly the
\textsf{weight}, \textsf{intercept}, and \textsf{noise} properties read at compile
time (\S\ref{sec:compile}). The snapshot reports itself parameterized only when
some value departs from the structural-only defaults, so a purely qualitative DAG
incurs no numeric machinery.

\paragraph{Interventions and counterfactuals}
The do-operator $do(X{=}x)$ fixes $X$'s mean to $x$, sets its variance to zero,
and drops its incoming edges; the general recursion then handles it with no
special case. Interventional means propagate in topological order,
$\mu(V)=\mu_{0,V}+\sum_P \beta_{P\to V}\,\mu(P)$; interventional variances follow
from an incremental covariance recursion in the same order; and the total effect
of $T$ on $Y$ is obtained by unit-impulse propagation
($\partial T{=}1$, $\partial V=\sum_P \beta_{P\to V}\,\partial P$, read off at
$Y$), avoiding any path enumeration. Counterfactuals follow the standard
abduction--action--prediction recipe over the same equations: infer the exogenous
noise consistent with the observed factual world, apply the hypothesized
intervention, and re-propagate. Every one of these reuses the topological order
already on the snapshot, so the quantitative layer adds semantics, not a second
data structure.

\paragraph{Worked example}
On a parameterized version of the running example, \lstinline[style=causaldsl]|INTERVENE ON interventions = 1 COMPARE TO interventions = 0 RETURN EFFECT ON reported_cases|
returns the total effect along the front-door paths, while a
\lstinline[style=causaldsl]|COUNTERFACTUAL| statement answers what
\textsf{reported\_cases} would have been for a specific observed unit under a
different intervention. Because $do(\cdot)$ yields a derived model---the DAG with
$X$'s in-edges removed---these statements register a new causal world in the
model registry (\S\ref{sec:lang:registry}) that subsequent structural queries can
address like any other, closing the loop between the qualitative and quantitative
layers over a single stored graph.

%% Faithful to LinearGaussianInference.java, CounterfactualInference.java,
%% DefaultCausalGraphRuntime.java. KEEP SHORT -- expressiveness, not a contribution.
%% ============================================================================
\section{Quantitative Layer and Expressiveness}\label{sec:quant}

The same compiled snapshot that answers structural queries also carries optional
numeric parameters, which let Causal4j answer \emph{quantitative} causal
questions---interventions and counterfactuals---without leaving the DBMS. We
describe the layer briefly: it is a demonstration of the expressiveness the
indexed engine enables, not a contribution of the paper.

\paragraph{Linear-Gaussian model}
When a model is parameterized, each variable follows the structural equation
$V = \mu_{0,V} + \sum_{P\in\mathrm{pa}(V)} \beta_{P\to V}\,P + \varepsilon_V$ with
$\varepsilon_V\sim\mathcal{N}(0,\sigma^2_V)$, where the edge coefficients
$\beta$, intercepts $\mu_0$, and noise variances $\sigma^2$ are exactly the
\textsf{weight}, \textsf{intercept}, and \textsf{noise} properties read at compile
time (\S\ref{sec:compile}). The snapshot reports itself parameterized only when
some value departs from the structural-only defaults, so a purely qualitative DAG
incurs no numeric machinery.

\paragraph{Interventions and counterfactuals}
The do-operator $do(X{=}x)$ fixes $X$'s mean to $x$, sets its variance to zero,
and drops its incoming edges; the general recursion then handles it with no
special case. Interventional means propagate in topological order,
$\mu(V)=\mu_{0,V}+\sum_P \beta_{P\to V}\,\mu(P)$; interventional variances follow
from an incremental covariance recursion in the same order; and the total effect
of $T$ on $Y$ is obtained by unit-impulse propagation
($\partial T{=}1$, $\partial V=\sum_P \beta_{P\to V}\,\partial P$, read off at
$Y$), avoiding any path enumeration. Counterfactuals follow the standard
abduction--action--prediction recipe over the same equations: infer the exogenous
noise consistent with the observed factual world, apply the hypothesized
intervention, and re-propagate. Every one of these reuses the topological order
already on the snapshot, so the quantitative layer adds semantics, not a second
data structure.

\paragraph{Worked example}
On a parameterized version of the running example, \lstinline[style=causaldsl]|INTERVENE ON interventions = 1 COMPARE TO interventions = 0 RETURN EFFECT ON reported_cases|
returns the total effect along the front-door paths, while a
\lstinline[style=causaldsl]|COUNTERFACTUAL| statement answers what
\textsf{reported\_cases} would have been for a specific observed unit under a
different intervention. Because $do(\cdot)$ yields a derived model---the DAG with
$X$'s in-edges removed---these statements register a new causal world in the
model registry (\S\ref{sec:lang:registry}) that subsequent structural queries can
address like any other, closing the loop between the qualitative and quantitative
layers over a single stored graph.

%% Faithful to LinearGaussianInference.java, CounterfactualInference.java,
%% DefaultCausalGraphRuntime.java. KEEP SHORT -- expressiveness, not a contribution.
%% ============================================================================
\section{Quantitative Layer and Expressiveness}\label{sec:quant}

The same compiled snapshot that answers structural queries also carries optional
numeric parameters, which let Causal4j answer \emph{quantitative} causal
questions---interventions and counterfactuals---without leaving the DBMS. We
describe the layer briefly: it is a demonstration of the expressiveness the
indexed engine enables, not a contribution of the paper.

\paragraph{Linear-Gaussian model}
When a model is parameterized, each variable follows the structural equation
$V = \mu_{0,V} + \sum_{P\in\mathrm{pa}(V)} \beta_{P\to V}\,P + \varepsilon_V$ with
$\varepsilon_V\sim\mathcal{N}(0,\sigma^2_V)$, where the edge coefficients
$\beta$, intercepts $\mu_0$, and noise variances $\sigma^2$ are exactly the
\textsf{weight}, \textsf{intercept}, and \textsf{noise} properties read at compile
time (\S\ref{sec:compile}). The snapshot reports itself parameterized only when
some value departs from the structural-only defaults, so a purely qualitative DAG
incurs no numeric machinery.

\paragraph{Interventions and counterfactuals}
The do-operator $do(X{=}x)$ fixes $X$'s mean to $x$, sets its variance to zero,
and drops its incoming edges; the general recursion then handles it with no
special case. Interventional means propagate in topological order,
$\mu(V)=\mu_{0,V}+\sum_P \beta_{P\to V}\,\mu(P)$; interventional variances follow
from an incremental covariance recursion in the same order; and the total effect
of $T$ on $Y$ is obtained by unit-impulse propagation
($\partial T{=}1$, $\partial V=\sum_P \beta_{P\to V}\,\partial P$, read off at
$Y$), avoiding any path enumeration. Counterfactuals follow the standard
abduction--action--prediction recipe over the same equations: infer the exogenous
noise consistent with the observed factual world, apply the hypothesized
intervention, and re-propagate. Every one of these reuses the topological order
already on the snapshot, so the quantitative layer adds semantics, not a second
data structure.

\paragraph{Worked example}
On a parameterized version of the running example, \lstinline[style=causaldsl]|INTERVENE ON interventions = 1 COMPARE TO interventions = 0 RETURN EFFECT ON reported_cases|
returns the total effect along the front-door paths, while a
\lstinline[style=causaldsl]|COUNTERFACTUAL| statement answers what
\textsf{reported\_cases} would have been for a specific observed unit under a
different intervention. Because $do(\cdot)$ yields a derived model---the DAG with
$X$'s in-edges removed---these statements register a new causal world in the
model registry (\S\ref{sec:lang:registry}) that subsequent structural queries can
address like any other, closing the loop between the qualitative and quantitative
layers over a single stored graph.
